\part{Scientific Positioning}

\chapter{Scientific Positioning and Comparative Analysis}
\label{chap:scientific_positioning}

\section{Synthesis of selected original research}

The original research papers reviewed in Chapter~\ref{chap:related_work} show that intelligent irrigation is not a single technology. It is a combination of field sensing, communication, data integration, prediction, simulation, recommendation, scheduling, and actuation. Low-cost IoT prototypes such as the work of Kodali and Sarjerao demonstrate the value of a simple sensor-to-pump loop using MQTT and a microcontroller \cite{kodali2017mqttirrigation}. Monitoring-oriented studies such as Balaceanu et al. show that practical irrigation decisions benefit from several agricultural variables rather than a single threshold \cite{balaceanu2019libelium}. IoT automation studies such as Gupta et al. connect sensing and control to water and energy efficiency objectives \cite{gupta2025iotirrigation}.

Machine-learning works add decision support to this sensing layer. Tancredi et al. connect sensor data, a Digital Twin direction, and machine-learning classification for irrigation management \cite{tancredi2025irrigationdt}. Ed-Daoudi et al. position predictive irrigation as a way to combine IoT data and adaptive scheduling \cite{eddaoudi2024predictiveirrigation}. Field Digital Twin studies go further by combining sensing, climate data, simulation, and irrigation scheduling. Bellvert et al. apply Digital Twin principles to automated vineyard irrigation scheduling \cite{bellvert2023vineyarddt}, while Millan et al. evaluate the Irri\_DesK Digital Twin in a commercial tomato farm \cite{millan2025tomatodt}.

Other works point toward future extensions. El Hachimi et al. combine Digital Twins and deep reinforcement learning for collective irrigation water distribution \cite{elhachimi2025collectiveirrigationdt}. Kallenberg et al. emphasize interoperability between agricultural Digital Twins, reinforcement-learning agents, and FIWARE-style context integration \cite{kallenberg2025interoperabledt}. Liu et al. focus on crop-level Digital Twin representation of cotton canopy development under water stress \cite{liu2025cottondt}. These works are useful for positioning, but they also show that advanced optimization, interoperability, and crop response modeling require validation conditions beyond a final-year prototype.

\section{Identified research and integration gaps}

The comparison of these papers reveals several gaps. First, monitoring systems often collect useful measurements but stop before prediction, simulation, recommendation, or physical control. Second, IoT control prototypes can switch pumps and demonstrate communication loops, but they usually do not maintain a synchronized virtual state or a traceable Digital Twin history. Third, machine-learning irrigation studies may report prediction metrics without integrating the model into a complete cyber-physical workflow from measurement to actuation. Fourth, field Digital Twin irrigation systems provide stronger agronomic value, but they require crop-specific calibration, field-scale deployment, and specialized data sources such as remote sensing or calibrated water-balance models. Fifth, reinforcement-learning and interoperability studies are important for future work but exceed the realistic scope of the implemented prototype.

The gap addressed by this thesis is therefore an integration gap. The objective is not to outperform field-validated Digital Twin irrigation systems. The objective is to assemble the main functions observed separately in the literature into a coherent supervised prototype: sensing, communication, storage, virtual state update, prediction, simulation, recommendation, validation, scheduling, actuation, and feedback.

\section{Positioning criteria}

The selected works are positioned using criteria derived from the functions required by an integrated Digital Twin irrigation workflow:

\begin{itemize}
    \item sensing: real agricultural or environmental measurements;
    \item communication: field-to-software data exchange;
    \item data integration: persistence and structured reuse of measurements;
    \item virtual state: a synchronized representation of the monitored zone;
    \item prediction: use of machine learning or predictive logic;
    \item simulation: evaluation of possible future states or irrigation scenarios;
    \item recommendation: generation of interpretable irrigation advice;
    \item scheduling: transformation of decisions into planned irrigation events;
    \item actuation: command of a pump, valve, or irrigation device;
    \item validation: prototype, dataset-level, simulation, or field-scale evidence;
    \item deployment feasibility: suitability for a constrained engineering prototype.
\end{itemize}

\section{Functional comparison of selected original works}

Table~\ref{tab:functional-comparison-original-works} compares the selected original works according to their main contribution and their limitation for this thesis. The purpose is not to rank the papers, but to clarify why the thesis is positioned as an integration-oriented prototype.

\begin{footnotesize}
\setlength{\tabcolsep}{2pt}
\renewcommand{\arraystretch}{1.15}
\begin{longtable}{@{}p{0.16\textwidth}p{0.18\textwidth}p{0.18\textwidth}p{0.18\textwidth}p{0.20\textwidth}@{}}
\caption{Functional comparison of selected original works}
\label{tab:functional-comparison-original-works}\\
\toprule
Study & Main focus & Strongest contribution & Digital Twin maturity & Limitation for this thesis \\
\midrule
\endfirsthead
\toprule
Study & Main focus & Strongest contribution & Digital Twin maturity & Limitation for this thesis \\
\midrule
\endhead
Kodali and Sarjerao \cite{kodali2017mqttirrigation} & Low-cost MQTT smart irrigation & Simple sensor-to-pump loop with remote status & Not a Digital Twin & Useful for MQTT/control background only \\
Balaceanu et al. \cite{balaceanu2019libelium} & Precision-agriculture monitoring & Multi-variable agricultural sensing & Monitoring layer only & No prediction, simulation, or actuation \\
Tancredi et al. \cite{tancredi2025irrigationdt} & ML-assisted irrigation decisions & Sensor-data classification for irrigation management & ML-enhanced Digital Twin direction & Limited field and implementation detail \\
Bellvert et al. \cite{bellvert2023vineyarddt} & Vineyard irrigation scheduling & Sentinel-2 assimilation and automated scheduling & Field-oriented irrigation Digital Twin & Requires remote-sensing and vineyard-specific calibration \\
Millan et al. \cite{millan2025tomatodt} & Tomato farm water-efficient management & Commercial-farm Digital Twin validation & Field-oriented irrigation Digital Twin & Site- and crop-specific management context \\
El Hachimi et al. \cite{elhachimi2025collectiveirrigationdt} & Collective irrigation-water distribution & Digital Twin with deep reinforcement learning & Optimization-oriented Digital Twin & Advanced optimization beyond prototype scope \\
Kallenberg et al. \cite{kallenberg2025interoperabledt} & Interoperable agricultural Digital Twins & FIWARE integration and RL recommendations & Interoperable AI-enhanced Digital Twin & General smart-farming scope, not irrigation-specific \\
Ed-Daoudi et al. \cite{eddaoudi2024predictiveirrigation} & Predictive smart irrigation & IoT and ML for adaptive irrigation scheduling & Predictive decision support & Limited explicit Digital Twin behavior \\
Gupta et al. \cite{gupta2025iotirrigation} & Automated IoT irrigation & Water and energy efficiency through automated control & IoT automation, not Digital Twin & No simulation or synchronized virtual state \\
Liu et al. \cite{liu2025cottondt} & Cotton water-stress Digital Twin & Crop-level water-stress representation & Crop Digital Twin & No direct irrigation-control workflow \\
\bottomrule
\end{longtable}
\end{footnotesize}

The table shows that no selected work covers all functions equally. IoT irrigation prototypes are practical and reproducible but are not full Digital Twins. Field Digital Twin systems are agronomically stronger but require calibration and validation conditions beyond the current project. Machine-learning studies support prediction but do not always close the cyber-physical loop. Reinforcement-learning and interoperability works are useful references for future evolution rather than direct implementation requirements.

\section{Positioning of the thesis contribution}

The thesis contribution is positioned between low-cost IoT irrigation prototypes and field-validated Digital Twin irrigation systems. Agrio is an integration-oriented prototype: it connects real sensing, MQTT communication, backend persistence, Digital Twin state update, Random Forest prediction, what-if simulation, recommendation, human validation, scheduling, and actuator control. This position is deliberately conservative. The system should not be described as a fully calibrated agronomic decision system or a production-grade irrigation optimizer.

Compared with low-cost IoT irrigation systems, the thesis adds structured persistence, Digital Twin state management, prediction, simulation, recommendation, and traceability. Compared with field Digital Twin systems, the thesis remains limited to prototype-scale validation and dataset-level ML evaluation. This makes the contribution useful as an engineering bridge between simple connected irrigation and more mature Digital Twin irrigation platforms.

\section{Explicit limitations}

The following limitations must remain explicit throughout the thesis:

\begin{itemize}
    \item no full-season field validation;
    \item no crop-specific calibration equivalent to commercial agronomic systems;
    \item no production-grade security claim;
    \item dataset-level ML validation only;
    \item prototype-scale actuation using a controlled environment;
    \item no implemented online learning, drift detection, crop-growth model, yield prediction, or full optimization engine.
\end{itemize}

\section{Summary}

This chapter merged the scientific positioning and comparative analysis into a single argument. The literature shows strong components in isolation, while the thesis addresses the engineering problem of connecting those components into a traceable Digital Twin-based irrigation prototype. The next chapter translates this positioning into requirements and design decisions.
